\chapter{Minimal Normal Subgroups and Transfer}

\section{Minimal Normal Subgroups}

Let $G$ be a finite group and let
\[
    1\neq N\lhd G.
\]
Then
\[
    \overline{G}=G/N
\]
is a group with
\[
    |\overline{G}|<|G|.
\]
Thus we have a short exact sequence
\[
    1\longrightarrow N\longrightarrow G\longrightarrow \overline{G}\longrightarrow 1.
\]
In this sense, $G$ is an extension of $N$ by $\overline{G}$.

\begin{lemma}[Structure of minimal normal subgroups]
Let $G$ be a finite group and let $N\lhd_{\mathrm{min}}G$ be a minimal normal subgroup of $G$. Then
\[
    N=S_1\times S_2\times\cdots\times S_r,
\]
where
\[
    S_1\cong S_2\cong\cdots\cong S_r
\]
are simple groups. Equivalently,
\[
    N\cong S^r
\]
for some simple group $S$.

In particular, either
\[
    N\cong C_p^r
\]
for some prime $p$, or
\[
    N\cong T^r
\]
for some nonabelian simple group $T$.
\end{lemma}

\begin{proof}
We argue by induction on $|N|$.

If $N$ is simple, then there is nothing to prove. So assume that $N$ is not simple. Choose a minimal normal subgroup
\[
    M\lhd_{\mathrm{min}}N.
\]
By induction applied to the group $N$ and its minimal normal subgroup $M$, we may write
\[
    M=T_1\times T_2\times\cdots\times T_m,
\]
where
\[
    T_1\cong T_2\cong\cdots\cong T_m
\]
are simple groups.

For every $g\in G$, since $N\lhd G$, the conjugate
\[
    M^g=g^{-1}Mg
\]
is again a minimal normal subgroup of $N$. Hence, by minimality,
\[
    M^g=M
    \qquad\text{or}\qquad
    M^g\cap M=1.
\]
More generally, any two distinct conjugates of $M$ under $G$ have trivial intersection.

Let
\[
    M^G=\gen{M^g:g\in G}
\]
be the normal closure of $M$ in $G$. Choose distinct conjugates
\[
    M^{g_1},M^{g_2},\ldots,M^{g_s}.
\]
Since these subgroups are normal in $N$, if $i\neq j$, then
\[
    [M^{g_i},M^{g_j}]
    \leq M^{g_i}\cap M^{g_j}
    =1.
\]
Therefore the distinct conjugates commute with one another, and so
\[
    M^G=M^{g_1}\times M^{g_2}\times\cdots\times M^{g_s}.
\]
Also,
\[
    1\neq M^G\lhd G
    \qquad\text{and}\qquad
    M^G\leq N.
\]
Since $N\lhd_{\mathrm{min}}G$, we must have
\[
    M^G=N.
\]
Thus $N$ is a direct product of simple groups, all isomorphic to the simple factors of $M$.

Finally, if the common simple factor is abelian, then it is cyclic of prime order, so
\[
    N\cong C_p^r.
\]
Otherwise the common simple factor is nonabelian simple.
\end{proof}

\begin{corollary}
If $G$ is solvable and $N\lhd_{\mathrm{min}}G$, then
\[
    N\cong C_p^r
\]
for some prime $p$ and some $r\geq1$.
\end{corollary}

\begin{proof}
Since $G$ is solvable, every subgroup of $G$ is solvable. Hence $N$ is solvable.

By the lemma,
\[
    N\cong S^r
\]
for some simple group $S$. Since $S$ is a quotient of $N$, the group $S$ is also solvable. A solvable simple group must be cyclic of prime order. Hence
\[
    S\cong C_p
\]
for some prime $p$, and therefore
\[
    N\cong C_p^r.
\]
\end{proof}

\section{Automorphisms of Nonabelian Simple Groups Products}

\begin{topic}[Question]
Let
\[
    M=T_1\times T_2\times\cdots\times T_r,
\]
where
\[
    T_1\cong T_2\cong\cdots\cong T_r\cong T
\]
and $T$ is nonabelian simple. We want to determine
\[
    \Aut(M).
\]
\end{topic}

\begin{lemma}
Let
\[
    M=T_1\times T_2\times\cdots\times T_r,
\]
where each $T_i$ is nonabelian simple. Then the minimal normal subgroups of $M$ are precisely
\[
    T_1,T_2,\ldots,T_r.
\]
\end{lemma}

\begin{proof}
Each $T_i$ is normal in $M$, and since $T_i$ is simple, it is a minimal normal subgroup of $M$.

Conversely, let $L\lhd_{\mathrm{min}}M$. Since $L\neq1$, there exists some $i$ such that
\[
    [L,T_i]\neq1.
\]
Indeed, if $[L,T_i]=1$ for every $i$, then $L$ centralizes all of $M$, so
\[
    L\leq Z(M)=1,
\]
because each $T_i$ is nonabelian simple and hence has trivial center. This is impossible.

Now
\[
    [L,T_i]\leq L\cap T_i.
\]
Thus
\[
    L\cap T_i\neq1.
\]
But $L\cap T_i\lhd M$ and $L$ is minimal normal in $M$, so
\[
    L\cap T_i=L.
\]
Hence $L\leq T_i$. Since $T_i$ is simple and $L\neq1$, we get
\[
    L=T_i.
\]
Therefore the minimal normal subgroups of $M$ are exactly the direct factors $T_i$.
\end{proof}

\begin{theorem}[Automorphisms of $T^r$]
Let $T$ be a nonabelian simple group and let
\[
    M=T_1\times T_2\times\cdots\times T_r,
    \qquad
    T_i\cong T.
\]
Then
\[
    \Aut(M)\cong \Aut(T)\wr S_r
    =
    \Aut(T)^r\rtimes S_r,
\]
where $S_r$ acts on $\Aut(T)^r$ by permuting the $r$ coordinates.

More explicitly, $\Aut(M)$ acts on the set
\[
    \{T_1,T_2,\ldots,T_r\},
\]
and the kernel of this action is
\[
    K=\Aut(T_1)\times\Aut(T_2)\times\cdots\times\Aut(T_r)
    \cong \Aut(T)^r.
\]
Moreover,
\[
    \Aut(M)/K\cong S_r.
\]
\end{theorem}

\begin{proof}
Let $\alpha\in\Aut(M)$. Since automorphisms preserve minimal normal subgroups, and since the minimal normal subgroups of $M$ are precisely
\[
    T_1,T_2,\ldots,T_r,
\]
the automorphism $\alpha$ permutes the direct factors. Thus we get a homomorphism
\[
    \Aut(M)\longrightarrow S_r.
\]

The kernel consists of those automorphisms that stabilize every factor $T_i$. If $\alpha$ lies in the kernel, then
\[
    \alpha|_{T_i}\in\Aut(T_i)
\]
for each $i$, and for
\[
    m=t_1t_2\cdots t_r
    \qquad(t_i\in T_i),
\]
we have
\[
    \alpha(m)
    =
    \alpha(t_1)\alpha(t_2)\cdots\alpha(t_r).
\]
Therefore
\[
    \Ker\left(\Aut(M)\to S_r\right)
    =
    \Aut(T_1)\times\cdots\times\Aut(T_r)
    =
    K.
\]

Conversely, every permutation of the factors gives an automorphism of $M$, since all the factors are isomorphic. Hence the map
\[
    \Aut(M)\longrightarrow S_r
\]
is surjective, and so
\[
    \Aut(M)/K\cong S_r.
\]

Thus
\[
    \Aut(M)\cong K\rtimes S_r
    \cong \Aut(T)^r\rtimes S_r
    =
    \Aut(T)\wr S_r.
\]
\end{proof}

\begin{example}[Minimal normal subgroups in groups of order $72$]
Let
\[
    |G|=72=2^3\cdot 3^2.
\]
Suppose that $G$ is solvable, and let
\[
    N\lhd_{\mathrm{min}}G.
\]
Then
\[
    N\cong C_p^r
\]
for some prime $p$.

Since $|N|\mid |G|=72$, the possibilities are
\[
    N\cong C_2,\ C_2^2,\ C_2^3,
    \qquad\text{or}\qquad
    N\cong C_3,\ C_3^2.
\]
Equivalently, after viewing $N$ as a vector space over $\mathbb F_p$, we have
\[
    N\leq C_2^3
    \qquad\text{or}\qquad
    N\leq C_3^2.
\]
\end{example}

\section{Burnside's Theorem and Transfer}

\begin{theorem}[Burnside]
Let $G$ be a finite group. If all Sylow subgroups of $G$ are cyclic, then $G$ is solvable.
\end{theorem}

\begin{remark}
A standard proof of this theorem uses the \textbf{transfer map}. We now define transfer and record its basic properties.
\end{remark}

\begin{definition}[Transfer]
Let $G$ be a finite group and let $H\leq G$. Put
\[
    \Omega=H\backslash G=\{Hx:x\in G\}.
\]
Choose a right transversal
\[
    \Omega=\{Hx_1,Hx_2,\ldots,Hx_n\},
    \qquad
    n=[G:H].
\]
For $g\in G$, right multiplication by $g$ gives a permutation of $\Omega$:
\[
    Hx_i\longmapsto Hx_i g.
\]
Thus there exists a permutation $\tau_g\in S_n$ such that
\[
    Hx_i g=Hx_{i^{\tau_g}}.
\]
Hence
\[
    x_i g x_{i^{\tau_g}}^{-1}\in H.
\]
Define
\[
    h_i(g)=x_i g x_{i^{\tau_g}}^{-1}\in H.
\]
The \textbf{transfer map} from $G$ to $H/H'$ is
\[
    V_{G\to H}:G\longrightarrow H/H',
\]
defined by
\[
    V_{G\to H}(g)
    =
    \prod_{i=1}^n h_i(g)H'
    =
    \prod_{i=1}^n
    \left(x_i g x_{i^{\tau_g}}^{-1}\right)H'.
\]
Here $H'=[H,H]$ is the commutator subgroup of $H$.
\end{definition}

\begin{lemma}[Basic properties of transfer]
Let $G$ be finite and let $H\leq G$.
\begin{enumerate}[label=(\arabic*)]
    \item The map
    \[
        V_{G\to H}:G\longrightarrow H/H'
    \]
    is a group homomorphism.

    \item The map $V_{G\to H}$ does not depend on the choice of right coset representatives
    \[
        x_1,x_2,\ldots,x_n.
    \]

    \item Let $g\in G$. Suppose the action of $\gen{g}$ on
    \[
        \Omega=H\backslash G
    \]
    has orbits
    \[
        (Hx_1,Hx_1g,\ldots,Hx_1g^{e_1-1}),
    \]
    \[
        (Hx_2,Hx_2g,\ldots,Hx_2g^{e_2-1}),
    \]
    \[
        \ldots,
    \]
    \[
        (Hx_t,Hx_tg,\ldots,Hx_tg^{e_t-1}).
    \]
    Then
    \[
        x_i g^{e_i}x_i^{-1}\in H
    \]
    for each $i$, and
    \[
        V_{G\to H}(g)
        =
        \prod_{i=1}^t
        \left(x_i g^{e_i}x_i^{-1}\right)H'.
    \]
\end{enumerate}
\end{lemma}

\begin{proof}
For each $g\in G$, write
\[
    Hx_i g=Hx_{i^{\tau_g}},
\]
and put
\[
    h_i(g)=x_i g x_{i^{\tau_g}}^{-1}\in H.
\]

\textbf{(1)} Let $g_1,g_2\in G$. Since right multiplication gives an action on $\Omega$, we have
\[
    i^{\tau_{g_1g_2}}
    =
    \left(i^{\tau_{g_1}}\right)^{\tau_{g_2}}.
\]
Hence
\[
\begin{aligned}
    h_i(g_1g_2)
    &=
    x_i g_1g_2 x_{i^{\tau_{g_1g_2}}}^{-1} \\[2mm]
    &=
    x_i g_1 x_{i^{\tau_{g_1}}}^{-1}
    \cdot
    x_{i^{\tau_{g_1}}}g_2x_{\left(i^{\tau_{g_1}}\right)^{\tau_{g_2}}}^{-1} \\[2mm]
    &=
    h_i(g_1)\,h_{i^{\tau_{g_1}}}(g_2).
\end{aligned}
\]
Therefore, in the abelian group $H/H'$,
\[
\begin{aligned}
    V_{G\to H}(g_1g_2)
    &=
    \prod_{i=1}^n h_i(g_1g_2)H' \\[2mm]
    &=
    \prod_{i=1}^n
    h_i(g_1)h_{i^{\tau_{g_1}}}(g_2)H' \\[2mm]
    &=
    \left(\prod_{i=1}^n h_i(g_1)H'\right)
    \left(\prod_{i=1}^n h_{i^{\tau_{g_1}}}(g_2)H'\right).
\end{aligned}
\]
Since $\tau_{g_1}$ is a permutation of $\{1,\ldots,n\}$,
\[
    \prod_{i=1}^n h_{i^{\tau_{g_1}}}(g_2)H'
    =
    \prod_{i=1}^n h_i(g_2)H'.
\]
Thus
\[
    V_{G\to H}(g_1g_2)
    =
    V_{G\to H}(g_1)V_{G\to H}(g_2).
\]
So $V_{G\to H}$ is a homomorphism.

\textbf{(2)} Let
\[
    y_1,y_2,\ldots,y_n
\]
be another choice of right coset representatives, with
\[
    Hy_i=Hx_i.
\]
Then there exist $t_i\in H$ such that
\[
    y_i=t_i x_i.
\]
For the new representatives, the corresponding element of $H$ is
\[
\begin{aligned}
    \widetilde{h}_i(g)
    &=
    y_i g y_{i^{\tau_g}}^{-1} \\[2mm]
    &=
    t_i x_i g x_{i^{\tau_g}}^{-1}t_{i^{\tau_g}}^{-1} \\[2mm]
    &=
    t_i h_i(g)t_{i^{\tau_g}}^{-1}.
\end{aligned}
\]
Hence in $H/H'$,
\[
\begin{aligned}
    \prod_{i=1}^n \widetilde{h}_i(g)H'
    &=
    \prod_{i=1}^n t_i h_i(g)t_{i^{\tau_g}}^{-1}H' \\[2mm]
    &=
    \left(\prod_{i=1}^n h_i(g)H'\right)
    \left(\prod_{i=1}^n t_iH'\right)
    \left(\prod_{i=1}^n t_{i^{\tau_g}}^{-1}H'\right).
\end{aligned}
\]
Since $\tau_g$ is a permutation,
\[
    \prod_{i=1}^n t_iH'
    =
    \prod_{i=1}^n t_{i^{\tau_g}}H'.
\]
Therefore
\[
    \prod_{i=1}^n \widetilde{h}_i(g)H'
    =
    \prod_{i=1}^n h_i(g)H',
\]
so the transfer does not depend on the chosen transversal.

\textbf{(3)} By part (2), we may choose a convenient transversal adapted to the $\gen{g}$-orbits. For an orbit
\[
    (Hx_i,Hx_ig,\ldots,Hx_ig^{e_i-1}),
\]
choose representatives
\[
    x_i,\ x_ig,\ \ldots,\ x_ig^{e_i-1}.
\]
For the first $e_i-1$ representatives in this orbit, the corresponding transfer terms are
\[
    (x_ig^k)g(x_ig^{k+1})^{-1}=1
    \qquad
    (0\leq k\leq e_i-2).
\]
For the last representative, the term is
\[
    (x_ig^{e_i-1})g x_i^{-1}
    =
    x_ig^{e_i}x_i^{-1}.
\]
Since
\[
    Hx_ig^{e_i}=Hx_i,
\]
we indeed have
\[
    x_ig^{e_i}x_i^{-1}\in H.
\]
Multiplying over all orbits gives
\[
    V_{G\to H}(g)
    =
    \prod_{i=1}^t
    \left(x_i g^{e_i}x_i^{-1}\right)H'.
\]
\end{proof}

\section{\texorpdfstring{$p$}{p}-Nilpotence and Burnside's Transfer Theorem}

\begin{definition}[$p$-nilpotent group]
Let $G$ be a finite group and let $p$ be a prime divisor of $|G|$. We say that $G$ is
\textbf{$p$-nilpotent} if
\[
    G=NP,
\]
where
\[
    N\lhd G,\qquad P\in \operatorname{Syl}_p(G),
    \qquad N\cap P=1.
\]
Equivalently, $G$ is $p$-nilpotent if and only if $G$ has a normal subgroup $N$ such that
\[
    |G/N|=|G|_p=|P|,
\]
where $P$ is a Sylow $p$-subgroup of $G$.
\end{definition}

\begin{theorem}[Burnside normal $p$-complement theorem]
Let $G$ be a finite group and let $P\in \operatorname{Syl}_p(G)$. If
\[
    N_G(P)=C_G(P),
\]
then $G$ is $p$-nilpotent.
\end{theorem}

\begin{proof}
We use transfer. Since
\[
    N_G(P)=C_G(P),
\]
and $P\leq N_G(P)$, we get
\[
    P\leq C_G(P).
\]
Hence $P$ is abelian, so
\[
    P'=1.
\]
Thus the transfer map gives a homomorphism
\[
    V_{G\to P}:G\longrightarrow P.
\]

Let $g\in P\setminus\{1\}$. Consider the action of $\gen{g}$ on
\[
    P\backslash G.
\]
Let the $\gen{g}$-orbits be represented by
\[
    Px_1,\ldots,Px_t,
\]
with orbit lengths
\[
    \ell_1,\ldots,\ell_t.
\]
By the orbit formula for transfer,
\[
    V_{G\to P}(g)
    =
    \prod_{i=1}^t x_i g^{\ell_i}x_i^{-1}.
\]
Here each factor lies in $P$:
\[
    x_i g^{\ell_i}x_i^{-1}\in P.
\]

We claim that
\[
    x_i g^{\ell_i}x_i^{-1}=g^{\ell_i}
\]
for every $i$.

Indeed, since $g^{\ell_i}\in P$ and $P$ is abelian, we have
\[
    P\leq C_G(g^{\ell_i}).
\]
Also, since
\[
    x_i g^{\ell_i}x_i^{-1}\in P
\]
and $P$ is abelian, we have
\[
    P\leq C_G(x_i g^{\ell_i}x_i^{-1}).
\]
But
\[
    C_G(x_i g^{\ell_i}x_i^{-1})
    =
    x_i C_G(g^{\ell_i})x_i^{-1}.
\]
Hence
\[
    x_i^{-1}Px_i\leq C_G(g^{\ell_i}).
\]
Therefore both
\[
    P
    \qquad\text{and}\qquad
    x_i^{-1}Px_i
\]
are Sylow $p$-subgroups of $C_G(g^{\ell_i})$. By Sylow's theorem, there exists
\[
    u_i\in C_G(g^{\ell_i})
\]
such that
\[
    u_i^{-1}x_i^{-1}Px_i u_i=P.
\]
Equivalently,
\[
    (x_i u_i)^{-1}P(x_i u_i)=P,
\]
so
\[
    x_i u_i\in N_G(P)=C_G(P).
\]
In particular, $x_i u_i$ centralizes $g^{\ell_i}$. Since $u_i\in C_G(g^{\ell_i})$, it follows that $x_i$ also centralizes $g^{\ell_i}$. Thus
\[
    x_i g^{\ell_i}x_i^{-1}=g^{\ell_i}.
\]

Therefore
\[
\begin{aligned}
    V_{G\to P}(g)
    &=
    \prod_{i=1}^t g^{\ell_i}  \\
    &=
    g^{\ell_1+\cdots+\ell_t}.
\end{aligned}
\]
Since the $\gen{g}$-orbits partition $P\backslash G$,
\[
    \ell_1+\cdots+\ell_t=[G:P].
\]
Hence
\[
    V_{G\to P}(g)=g^{[G:P]}.
\]
Because $g$ has $p$-power order and $p\nmid [G:P]$, we have
\[
    g^{[G:P]}\neq 1.
\]
Thus
\[
    V_{G\to P}|_P:P\longrightarrow P
\]
is injective. Since $P$ is finite, it is also surjective. Hence
\[
    V_{G\to P}(P)=P.
\]
Therefore
\[
    V_{G\to P}(G)=P.
\]

Let
\[
    K=\Ker V_{G\to P}.
\]
Then
\[
    K\lhd G,
\]
and by the first isomorphism theorem,
\[
    G/K\cong P.
\]
Hence
\[
    |K|=\frac{|G|}{|P|}.
\]
Moreover, since $V_{G\to P}|_P$ is injective,
\[
    K\cap P=1.
\]
Thus
\[
    |KP|=\frac{|K||P|}{|K\cap P|}
    =
    |G|,
\]
so
\[
    G=KP.
\]
Therefore $G$ is $p$-nilpotent.
\end{proof}

\begin{theorem}[Burnside]
Let $G$ be a finite group whose Sylow subgroups are all cyclic. Let $p$ be the smallest prime divisor of $|G|$, and let
\[
    P\in \operatorname{Syl}_p(G).
\]
Then:
\begin{enumerate}[label=(\arabic*)]
    \item $G$ is $p$-nilpotent.
    \item $G$ is solvable.
\end{enumerate}
\end{theorem}

\begin{proof}
Since $P$ is cyclic, write
\[
    |P|=p^e.
\]
Then
\[
    |\Aut(P)|=p^{e-1}(p-1).
\]
Hence every prime divisor of $|\Aut(P)|$ is at most $p$.

Conjugation induces an embedding
\[
    N_G(P)/C_G(P)\hookrightarrow \Aut(P).
\]
Thus every prime divisor of
\[
    |N_G(P)/C_G(P)|
\]
is at most $p$.

On the other hand, since $P$ is cyclic, it is abelian. Hence
\[
    P\leq C_G(P).
\]
Therefore
\[
    |N_G(P)/C_G(P)|
    \mid
    |N_G(P)/P|.
\]
Also
\[
    |N_G(P)/P|\mid [G:P].
\]
Since $p$ is the smallest prime divisor of $|G|$ and $p\nmid [G:P]$, every prime divisor of $[G:P]$ is strictly larger than $p$.

Thus every prime divisor of $|N_G(P)/C_G(P)|$ is both at most $p$ and strictly larger than $p$, which is impossible unless
\[
    |N_G(P)/C_G(P)|=1.
\]
Hence
\[
    N_G(P)=C_G(P).
\]
By Burnside's normal $p$-complement theorem, $G$ is $p$-nilpotent.

So there exists a normal subgroup
\[
    K\lhd G
\]
such that
\[
    |K|=\frac{|G|}{|P|}.
\]
Every Sylow subgroup of $K$ is also a Sylow subgroup of $G$ for the corresponding prime, and hence is cyclic.

Now let $p_1$ be the smallest prime divisor of $|K|$. By the same argument, $K$ is $p_1$-nilpotent. Hence $K$ has a normal subgroup
\[
    K_1\lhd K
\]
such that
\[
    |K_1|=\frac{|K|}{|P_1|},
\]
where
\[
    P_1\in \operatorname{Syl}_{p_1}(K).
\]
Repeating this process, we obtain a normal series
\[
    G=K_0\triangleright K_1\triangleright K_2\triangleright\cdots\triangleright K_m=Q\triangleright 1,
\]
where $Q$ is a Sylow $q$-subgroup of $G$ and $q$ is the largest prime divisor of $|G|$.

Each factor
\[
    K_i/K_{i+1}
\]
is cyclic, and $Q$ is cyclic. Hence this is a normal series with abelian factors. Therefore $G$ is solvable.
\end{proof}

\begin{remark}
In the above process, the final Sylow subgroup
\[
    Q\in \operatorname{Syl}_q(G)
\]
is actually normal in $G$.
\end{remark}

\begin{proof}
Suppose that
\[
    Q\lhd K_n\lhd K_{n-1}.
\]
For every $g\in K_{n-1}$, since $K_n\lhd K_{n-1}$, we have
\[
    g^{-1}Qg\leq g^{-1}K_n g=K_n.
\]
Thus $g^{-1}Qg$ is a Sylow $q$-subgroup of $K_n$. Since
\[
    Q\lhd K_n,
\]
the group $K_n$ has a unique Sylow $q$-subgroup. Hence
\[
    g^{-1}Qg=Q.
\]
Therefore
\[
    Q\lhd K_{n-1}.
\]
Repeating this argument upward along the series gives
\[
    Q\lhd G.
\]
\end{proof}

\begin{corollary}
If a Sylow $2$-subgroup of $G$ is cyclic, then $G$ is $2$-nilpotent. In particular, $G$ is solvable.
\end{corollary}

\begin{proof}
Let
\[
    P\in \operatorname{Syl}_2(G)
\]
be cyclic.

If $|G|$ is odd, then $G$ is solvable by the Feit--Thompson theorem. So assume that $2\mid |G|$.

Since $P$ is cyclic, it is abelian, so
\[
    P\leq C_G(P).
\]
Again conjugation induces
\[
    N_G(P)/C_G(P)\hookrightarrow \Aut(P).
\]
But $\Aut(P)$ is a $2$-group. Hence
\[
    N_G(P)/C_G(P)
\]
is a $2$-group.

On the other hand,
\[
    |N_G(P)/C_G(P)|
    \mid
    |N_G(P)/P|
    \mid
    [G:P],
\]
and $[G:P]$ is odd. Thus
\[
    N_G(P)/C_G(P)=1.
\]
Therefore
\[
    N_G(P)=C_G(P).
\]
By Burnside's normal $p$-complement theorem, $G$ is $2$-nilpotent. Hence there exists
\[
    N\lhd G
\]
such that
\[
    G=NP,
    \qquad
    N\cap P=1.
\]
Then $|N|$ is odd, so $N$ is solvable by the Feit--Thompson theorem. Also
\[
    G/N\cong P
\]
is cyclic, hence solvable. Since both $N$ and $G/N$ are solvable, $G$ is solvable.
\end{proof}

\section{Consequences of Burnside's Normal Complement Theorem}

\begin{proposition}
Let $G$ be a finite nonabelian simple group, and let $p$ be a prime divisor of $|G|$.
\begin{enumerate}[label=(\arabic*)]
    \item If $p$ is the smallest prime divisor of $|G|$, then
    \[
        p^3\mid |G|
        \qquad\text{or}\qquad
        12\mid |G|.
    \]
    \item If
    \[
        |G|=pm,
        \qquad
        (p,m)=1,
    \]
    and $P\in \operatorname{Syl}_p(G)$, then
    \[
        C_G(P)\neq N_G(P)\neq G,
    \]
    and
    \[
        \left|N_G(P)/C_G(P)\right|\mid (p-1).
    \]
\end{enumerate}
\end{proposition}

\begin{proof}
We first prove (2). Since $|G|=pm$ with $(p,m)=1$, a Sylow $p$-subgroup has order $p$. Hence
\[
    P\cong C_p.
\]
Conjugation gives a homomorphism
\[
    N_G(P)\longrightarrow \Aut(P),
\]
whose kernel is $C_G(P)$. Therefore
\[
    N_G(P)/C_G(P)\hookrightarrow \Aut(P).
\]
Since
\[
    \Aut(C_p)\cong C_{p-1},
\]
we get
\[
    \left|N_G(P)/C_G(P)\right|\mid (p-1).
\]

Also $N_G(P)\neq G$, because otherwise $P\lhd G$, contradicting the simplicity of $G$.

If $N_G(P)=C_G(P)$, then by Burnside's normal $p$-complement theorem, $G$ is $p$-nilpotent. Thus $G$ has a normal $p$-complement, say $K\lhd G$, with
\[
    |K|=\frac{|G|}{|P|}=m.
\]
Since $G$ is nonabelian simple, this is impossible. Hence
\[
    C_G(P)\neq N_G(P).
\]
This proves (2).

Now prove (1). Assume that $p$ is the smallest prime divisor of $|G|$, and suppose that
\[
    p^3\nmid |G|.
\]
We shall prove that $12\mid |G|$.

First, $p^2$ must divide $|G|$. Indeed, if $p^2\nmid |G|$, then $|G|=pm$ with $(p,m)=1$. By part (2),
\[
    C_G(P)\neq N_G(P),
\]
and
\[
    \left|N_G(P)/C_G(P)\right|\mid (p-1).
\]
But $P\leq C_G(P)$, so
\[
    \left|N_G(P)/C_G(P)\right|
    \mid
    \left|N_G(P)/P\right|
    \mid
    \frac{|G|}{p}.
\]
Since $p$ is the smallest prime divisor of $|G|$ and $p\nmid |G|/p$, every prime divisor of $|G|/p$ is strictly larger than $p$. This contradicts the fact that
\[
    \left|N_G(P)/C_G(P)\right|
\]
is a nontrivial divisor of $p-1$. Therefore
\[
    p^2\mid |G|.
\]

Since we assumed $p^3\nmid |G|$, we have
\[
    p^2\parallel |G|.
\]
Let
\[
    P\in \operatorname{Syl}_p(G).
\]
Then
\[
    |P|=p^2.
\]
Hence either
\[
    P\cong C_{p^2}
    \qquad\text{or}\qquad
    P\cong C_p\times C_p.
\]

Suppose first that
\[
    P\cong C_{p^2}.
\]
Then
\[
    |\Aut(P)|=\varphi(p^2)=p(p-1).
\]
Again,
\[
    N_G(P)/C_G(P)\hookrightarrow \Aut(P).
\]
Since $P$ is cyclic, it is abelian, so
\[
    P\leq C_G(P).
\]
Thus
\[
    \left|N_G(P)/C_G(P)\right|
    \mid
    \frac{|G|}{p^2}.
\]
Every prime divisor of $|G|/p^2$ is strictly larger than $p$, whereas every prime divisor of $p(p-1)$ is at most $p$. Hence
\[
    N_G(P)=C_G(P).
\]
By Burnside's normal $p$-complement theorem, $G$ is $p$-nilpotent, contradicting the simplicity of $G$. Therefore
\[
    P\not\cong C_{p^2}.
\]
So
\[
    P\cong C_p\times C_p.
\]

Now $P$ is abelian, hence
\[
    P\leq C_G(P).
\]
If
\[
    N_G(P)=C_G(P),
\]
then again Burnside's normal $p$-complement theorem implies that $G$ is $p$-nilpotent, impossible. Thus
\[
    N_G(P)\neq C_G(P).
\]
Hence the quotient
\[
    N_G(P)/C_G(P)
\]
is nontrivial.

Since
\[
    P\cong C_p\times C_p,
\]
we have
\[
    \Aut(P)\cong \operatorname{GL}(2,p).
\]
Thus
\[
    N_G(P)/C_G(P)\hookrightarrow \operatorname{GL}(2,p).
\]
Moreover,
\[
    |\operatorname{GL}(2,p)|
    =
    (p^2-1)(p^2-p)
    =
    p(p-1)^2(p+1).
\]
Let $r$ be a prime divisor of
\[
    \left|N_G(P)/C_G(P)\right|.
\]
Since $P\leq C_G(P)$, we have
\[
    \left|N_G(P)/C_G(P)\right|
    \mid
    \frac{|G|}{p^2}.
\]
Thus $r$ is a prime divisor of $|G|/p^2$. Since $p$ is the smallest prime divisor of $|G|$ and $p^3\nmid |G|$, we get
\[
    r>p.
\]
On the other hand, $r$ divides
\[
    |\operatorname{GL}(2,p)|=p(p-1)^2(p+1).
\]
Since $r>p$, it cannot divide $p$ or $p-1$. Therefore
\[
    r\mid p+1.
\]
Thus
\[
    p<r\leq p+1.
\]
Hence
\[
    r=p+1.
\]
Since $r$ is prime, this forces
\[
    p=2,
    \qquad
    r=3.
\]
Therefore
\[
    |P|=p^2=4
\]
and $3\mid |G|/4$. Hence
\[
    12\mid |G|.
\]
This proves (1).
\end{proof}

\begin{example}
If
\[
    |G|=5\cdot 7\cdot 13,
\]
then $G$ is cyclic.

Let $G_p$ denote a Sylow $p$-subgroup of $G$.

First consider $p=13$. Let
\[
    G_{13}\in \operatorname{Syl}_{13}(G).
\]
The number $n_{13}$ of Sylow $13$-subgroups satisfies
\[
    n_{13}\equiv 1\pmod{13},
    \qquad
    n_{13}\mid 35.
\]
The divisors of $35$ are $1,5,7,35$, and only $1$ is congruent to $1$ modulo $13$. Hence
\[
    n_{13}=1.
\]
Thus
\[
    G_{13}\lhd G.
\]

Now conjugation gives
\[
    G/C_G(G_{13})
    =
    N_G(G_{13})/C_G(G_{13})
    \hookrightarrow
    \Aut(G_{13}).
\]
Since
\[
    G_{13}\cong C_{13},
\]
we have
\[
    \Aut(G_{13})\cong C_{12}.
\]
Also $G_{13}\leq C_G(G_{13})$, so
\[
    |G/C_G(G_{13})|
    \mid
    \frac{|G|}{13}
    =
    35.
\]
Therefore
\[
    |G/C_G(G_{13})|
    \mid \gcd(35,12)=1.
\]
Hence
\[
    C_G(G_{13})=G.
\]
So
\[
    G_{13}\leq Z(G).
\]

By Burnside's normal $p$-complement theorem applied to $p=13$, the group $G$ is $13$-nilpotent. Hence there exists a normal subgroup
\[
    K\lhd G
\]
such that
\[
    |K|=5\cdot 7=35
\]
and
\[
    G=K G_{13},
    \qquad
    K\cap G_{13}=1.
\]
Since $G_{13}\leq Z(G)$, this product is direct:
\[
    G=K\times G_{13}.
\]

It remains to show that $K$ is cyclic. Since
\[
    |K|=35=5\cdot 7,
\]
let
\[
    Q\in \operatorname{Syl}_7(K).
\]
The number $n_7(K)$ of Sylow $7$-subgroups of $K$ satisfies
\[
    n_7(K)\equiv 1\pmod 7,
    \qquad
    n_7(K)\mid 5.
\]
Hence
\[
    n_7(K)=1.
\]
Thus
\[
    Q\lhd K.
\]
Conjugation gives
\[
    K/C_K(Q)
    =
    N_K(Q)/C_K(Q)
    \hookrightarrow
    \Aut(Q)
    \cong C_6.
\]
Since $Q\leq C_K(Q)$, the order of $K/C_K(Q)$ divides
\[
    |K/Q|=5.
\]
Thus
\[
    |K/C_K(Q)|\mid \gcd(5,6)=1,
\]
so
\[
    C_K(Q)=K.
\]
Therefore
\[
    Q\leq Z(K).
\]

Let
\[
    R\in \operatorname{Syl}_5(K).
\]
Then $R$ centralizes $Q$, and
\[
    K=QR,
    \qquad
    Q\cap R=1.
\]
Thus
\[
    K\cong Q\times R\cong C_7\times C_5\cong C_{35}.
\]
Consequently
\[
    G\cong K\times G_{13}
    \cong C_{35}\times C_{13}
    \cong C_{455}.
\]
Hence $G$ is cyclic.
\end{example}